\chapter{Inventaire des controllers backend}

\begingroup
\small
\setlength{\tabcolsep}{4pt}
\begin{longtable}{>{\raggedright\arraybackslash}p{4.3cm} >{\raggedright\arraybackslash}p{9cm}}
\toprule
Controller & Responsabilite principale \\
\midrule
\code{AgentWorkloadController} & Charge de travail des agents et vues de supervision. \\
\code{AIAssistantController} & Endpoints d'analyse, diagnostic, suggestion, copilote et tendances. \\
\code{AttachmentController} & Upload, recuperation et suppression des pieces jointes. \\
\code{AuthController} & Authentification locale ou aides de developpement selon les profils. \\
\code{CamundaController} & Consultation ou pilotage d'informations liees au moteur BPM. \\
\code{ClientController} & CRUD et recherche sur les clients. \\
\code{CommentController} & Gestion des commentaires publics et internes. \\
\code{DashboardController} & KPIs et statistiques. \\
\code{EscalationPolicyController} & Regles et configuration d'escalade. \\
\code{KnowledgeBaseController} & Base de connaissances. \\
\code{NotificationController} & Notifications utilisateur. \\
\code{ReportController} & Export et rapports. \\
\code{SatisfactionController} & Sondages et indicateurs de satisfaction. \\
\code{SupportCategoryController} & Categorie de support et referentiel associe. \\
\code{TicketController} & Coeur des actions ticket. \\
\code{UserController} & Gestion des utilisateurs. \\
\bottomrule
\end{longtable}
\endgroup

\chapter{Inventaire des services backend}

\begingroup
\small
\setlength{\tabcolsep}{4pt}
\begin{longtable}{>{\raggedright\arraybackslash}p{4.5cm} >{\raggedright\arraybackslash}p{8.8cm}}
\toprule
Service & Responsabilite principale \\
\midrule
\code{AgentSkillService} & Gestion des competences agents. \\
\code{AgentWorkloadService} & Calcul de charge et supervision des agents. \\
\code{AlfrescoArchiveService} & Archivage metier vers la GED. \\
\code{AlfrescoCmisService} & Communication CMIS avec Alfresco. \\
\code{AttachmentService} & Gestion des fichiers joints. \\
\code{BusinessHoursService} & Regles d'horaires ouvrables et calculs associes. \\
\code{CamundaAsyncService} & Traitements Camunda asynchrones. \\
\code{CamundaBootstrapService} & Resynchronisation des processus au demarrage. \\
\code{CamundaService} & Integration principale avec le moteur BPMN. \\
\code{ClientService} & Logique metier client. \\
\code{CommentService} & Commentaires et regles associees. \\
\code{DashboardService} & KPI et statistiques. \\
\code{EscalationService} & Escalade et shortlist d'assignation. \\
\code{KeycloakAdminService} & Interactions d'administration avec Keycloak. \\
\code{KnowledgeBaseService} & Capitalisation des solutions. \\
\code{NotificationService} & Generation et diffusion des notifications. \\
\code{ReportService} & Rapports PDF et Excel. \\
\code{SatisfactionService} & Gestion de la satisfaction client. \\
\code{SlaTimerService} & Reaction aux seuils de SLA. \\
\code{SupportCategoryService} & Gestion des categories de support. \\
\code{TicketAutomationService} & Automatisation de taches autour du ticket. \\
\code{TicketService} & Cycle de vie principal des tickets. \\
\code{UserIdentityService} & Resolution de l'identite courante. \\
\code{UserService} & Gestion metier des utilisateurs. \\
\bottomrule
\end{longtable}
\endgroup

\chapter{Inventaire des services Docker}

\begingroup
\small
\setlength{\tabcolsep}{4pt}
\begin{longtable}{>{\raggedright\arraybackslash}p{2.7cm} >{\raggedright\arraybackslash}p{2.2cm} >{\raggedright\arraybackslash}p{1.8cm} >{\raggedright\arraybackslash}p{7cm}}
\toprule
Service & Profil & Port & Observation \\
\midrule
\code{mysql} & standard & 3306 & Base principale avec healthcheck. \\
\code{keycloak} & standard & 8180 & IAM du systeme. \\
\code{backend} & standard & 8082 & API et Camunda integres. \\
\code{frontend} & standard & 4200 & UI Angular servie par Nginx. \\
\code{alfresco-postgres} & standard & interne & DB du referentiel documentaire. \\
\code{alfresco} & standard & 8090 & Repository documentaire. \\
\code{share} & standard & 8091 & UI Alfresco Share. \\
\code{sonarqube} & tools & 9000 & Analyse qualite optionnelle. \\
\code{mailhog} & tools & 8025/1025 & Serveur SMTP de test. \\
\code{ollama} & standard & 11434 & Runtime des modeles locaux. \\
\code{ollama-init} & standard & interne & Pull du modele au demarrage. \\
\code{ai-agent} & standard & 8000 & FastAPI pour les fonctions IA. \\
\code{demo-data-loader} & demo & interne & Rechargement controle du seed de demo. \\
\bottomrule
\end{longtable}
\endgroup

\chapter{Inventaire des modules frontend}

\begingroup
\small
\setlength{\tabcolsep}{4pt}
\begin{longtable}{>{\raggedright\arraybackslash}p{3.8cm} >{\raggedright\arraybackslash}p{9.2cm}}
\toprule
Module frontend & Description \\
\midrule
\code{dashboard} & Vue de pilotage avec KPI, tendances, repartitions et performance. \\
\code{tickets} & Liste, detail, creation, edition, assignation, escalade, workflow et historique. \\
\code{clients} & Gestion des clients, recherche, detail et actions associees. \\
\code{users} & Gestion des utilisateurs et de leurs profils. \\
\code{profile} & Espace personnel de l'utilisateur connecte. \\
\code{archives-reports} & Consultation et exploitation des archives et rapports. \\
\code{ai-assistant} & Chat IA, diagnostic, suggestion, analyse ticket et tendances. \\
\code{layout/header} & Barre haute, badges de notifications, menu utilisateur. \\
\code{layout/sidebar} & Navigation laterale adaptee aux roles. \\
\code{core/services} & Services HTTP et fonctions transverses. \\
\bottomrule
\end{longtable}
\endgroup

\chapter{Matrice simplifiee des permissions}

\begingroup
\footnotesize
\setlength{\tabcolsep}{4pt}
\begin{longtable}{>{\raggedright\arraybackslash}p{4.6cm} >{\centering\arraybackslash}p{1.8cm} >{\centering\arraybackslash}p{2cm} >{\centering\arraybackslash}p{2cm} >{\centering\arraybackslash}p{2cm}}
\toprule
Action & Client & Agent & Manager & Admin \\
\midrule
Creer un ticket & Oui & Oui & Oui & Oui \\
Voir ses tickets & Oui & N/A & N/A & N/A \\
Voir tickets assignes & Non & Oui & Oui & Oui \\
Voir tous les tickets & Non & Non & Oui & Oui \\
Assigner un ticket & Non & Non & Oui & Oui \\
Prendre le ticket & Non & Oui si autorise & Oui & Oui \\
Resoudre un ticket & Non & Oui si assigne & Oui selon contexte & Oui selon contexte \\
Escalader & Non & Oui si assigne & Oui & Oui \\
Gerer le SLA & Non & Oui si assigne & Oui & Oui \\
Demander une revue manager & Non & Non & Oui & Oui \\
Cloturer & Oui selon parcours & Non & Oui & Oui \\
Archiver & Non & Non & Oui selon regle & Oui \\
Voir commentaires internes & Non & Oui & Oui & Oui \\
Gerer utilisateurs & Non & Non & Oui selon ecrans & Oui \\
Gerer clients & Non & Non & Oui & Oui \\
\bottomrule
\end{longtable}
\endgroup

\chapter{Catalogue fonctionnel des endpoints metiers}

\begingroup
\footnotesize
\setlength{\tabcolsep}{4pt}
\begin{longtable}{>{\raggedright\arraybackslash}p{1.7cm} >{\raggedright\arraybackslash}p{5.5cm} >{\raggedright\arraybackslash}p{5.5cm}}
\toprule
Methode & Endpoint & Usage \\
\midrule
GET & \code{/api/tickets} & Liste paginee des tickets. \\
GET & \code{/api/tickets/search} & Recherche multicritere. \\
GET & \code{/api/tickets/\{id\}} & Detail d'un ticket. \\
GET & \code{/api/tickets/reference/\{ref\}} & Recuperation par reference. \\
POST & \code{/api/tickets} & Creation d'un ticket. \\
PUT & \code{/api/tickets/\{id\}} & Mise a jour d'un ticket. \\
POST & \code{/api/tickets/\{id\}/assign/\{agentId\}} & Assignation. \\
POST & \code{/api/tickets/\{id\}/take-charge} & Prise en charge. \\
POST & \code{/api/tickets/\{id\}/resolve} & Resolution. \\
POST & \code{/api/tickets/\{id\}/close} & Cloture. \\
POST & \code{/api/tickets/\{id\}/archive} & Archivage. \\
POST & \code{/api/tickets/\{id\}/escalate} & Escalade manuelle. \\
POST & \code{/api/tickets/\{id\}/manager-review} & Revue manager. \\
POST & \code{/api/tickets/\{id\}/sla-pause} & Pause SLA. \\
POST & \code{/api/tickets/\{id\}/sla-resume} & Reprise SLA. \\
POST & \code{/api/tickets/\{id\}/sla-extend} & Extension SLA. \\
POST & \code{/api/tickets/\{id\}/reject-resolution} & Refus de resolution. \\
GET & \code{/api/tickets/\{id\}/comments} & Commentaires complets. \\
GET & \code{/api/tickets/\{id\}/comments/public} & Commentaires publics. \\
POST & \code{/api/tickets/\{id\}/comments} & Ajout de commentaire. \\
GET & \code{/api/tickets/\{id\}/attachments} & Liste des pieces jointes. \\
POST & \code{/api/tickets/\{id\}/attachments} & Upload d'une piece jointe. \\
GET & \code{/api/tickets/\{id\}/history} & Historique du ticket. \\
GET & \code{/api/tickets/\{id\}/workflow-status} & Etat du workflow. \\
GET & \code{/api/tickets/\{id\}/workflow-trace} & Trace du workflow. \\
GET & \code{/api/clients} & Liste des clients. \\
GET & \code{/api/users} & Liste des utilisateurs. \\
GET & \code{/api/dashboard/*} & KPI et statistiques. \\
GET & \code{/api/notifications/*} & Recuperation des notifications. \\
GET & \code{/api/ai/status} & Etat de la brique IA. \\
GET & \code{/api/ai/analyze/\{id\}} & Analyse IA d'un ticket. \\
GET & \code{/api/ai/diagnose/\{id\}} & Diagnostic IA. \\
GET & \code{/api/ai/suggest-response/\{id\}} & Suggestion de reponse. \\
GET & \code{/api/ai/escalation-summary/\{id\}} & Resume d'escalade. \\
POST & \code{/api/ai/chat} & Chat contextualise. \\
\bottomrule
\end{longtable}
\endgroup

\chapter{Scenarios de tests et d'acceptation}

\begin{longtable}{>{\raggedright\arraybackslash}p{1.2cm} >{\raggedright\arraybackslash}p{4.2cm} >{\raggedright\arraybackslash}p{7.6cm}}
\toprule
ID & Scenario & Resultat attendu \\
\midrule
T1 & Creation d'un ticket par un client & Le ticket est cree, reference generee, workflow lance. \\
T2 & Creation d'un ticket sans client cote staff & Le systeme bloque si la regle client obligatoire s'applique. \\
T3 & Assignation par manager & Le ticket change d'etat et l'agent est notifie. \\
T4 & Assignation d'un ticket finalise & Action refusee. \\
T5 & Prise en charge par l'agent assigne & Le ticket passe dans un etat compatible. \\
T6 & Resolution par un agent non assigne & Action refusee. \\
T7 & Escalade manuelle par agent assigne & Historique et destinataire d'escalade mis a jour. \\
T8 & Revue manager sur ticket non finalise & Action autorisee aux bons profils. \\
T9 & Pause SLA sur ticket actif & Etat SLA passe en pause avec tracabilite. \\
T10 & Reprise SLA & Chronometrage repris correctement. \\
T11 & Ticket a risque a 75\% & Alerte visible cote interface et notification envoyee. \\
T12 & Ticket depasse a 100\% & Processus d'escalade ou supervision declenche. \\
T13 & Validation client de la resolution & Workflow progresse vers archivage. \\
T14 & Refus client de la resolution & Retour en traitement ou en qualification selon le flux. \\
T15 & Consultation Camunda Cockpit & Les instances visibles correspondent aux tickets actifs. \\
T16 & Upload de piece jointe & Le fichier apparait dans le ticket. \\
T17 & Archivage documentaire & Rapport et dossier archives accessibles. \\
T18 & Consultation des notifications temps reel & Les badges et compteurs se mettent a jour. \\
T19 & Utilisation de l'assistant IA avec reference ticket & La reference est resolue sans erreur technique. \\
T20 & Tri des clients par raison sociale & La liste charge sans erreur serveur. \\
T21 & Chargement de la modale d'assignation & Une liste s'affiche meme si l'IA est lente. \\
T22 & Affichage des candidats non eligibles & Les badges expliquent pourquoi certains agents ne sont pas selectionnables. \\
T23 & Rechargement du seed Docker & La base et Camunda restent coherents. \\
T24 & Connexion manager & Les boutons manager sont visibles et fonctionnels. \\
T25 & Connexion agent & Les actions de supervision non autorisees ne sont pas proposees. \\
T26 & Connexion client & Le client ne voit que son perimetre. \\
T27 & Build backend & La compilation Maven doit aboutir. \\
T28 & Build frontend & Le build Angular doit aboutir. \\
T29 & Docker healthchecks & Les conteneurs critiques doivent passer en healthy. \\
T30 & Parcours complet incident critique & De la creation a l'archivage, toutes les etapes restent coherentes. \\
\bottomrule
\end{longtable}

\chapter{Journal detaille de corrections recentes}

\section{Creation de ticket plus rapide}

Le chemin de creation du ticket a ete revu pour eviter que des traitements secondaires retardent trop la reponse visible par l'utilisateur. Cette correction ameliore nettement la sensation de fluidite.

\section{Resynchronisation Camunda}

L'initialisation de workflows pour des donnees semees directement en base a ete prise en charge afin de rendre le moteur BPM de nouveau coherent avec les tickets existants.

\section{Clarification des autorisations}

Les autorisations ont ete revues pour que le frontend et le backend racontent la meme chose. Cette clarification ameliore la credibilite de l'outil.

\section{Modale d'assignation plus robuste}

La modale charge maintenant une liste de secours avant de tenter un enrichissement IA, ce qui evite les ecrans de chargement infinis.

\section{Affichage des candidats non eligibles}

L'utilisateur voit desormais pourquoi certains profils ne sont pas selectionnables. Cette transparence est precieuse pour la comprehension metier.

\section{Resolution ID ou reference dans l'outil IA}

L'assistant IA accepte desormais un identifiant interne ou une reference visible comme \code{SF-1013}. Cela rapproche l'outil du langage reel des utilisateurs.

\section{Correction du tri des clients}

Le frontend envoie maintenant un champ de tri compatible avec le backend, ce qui evite une erreur serveur evitable.

\section{Amelioration de la visibilite des alertes}

Des compteurs et badges relies a de vraies donnees ont ete ajoutes pour mettre en evidence les actions requises et les alertes sur les tickets.

\section{Jeu de donnees de demonstration plus realiste}

La base a ete rechargee avec des donnees plus riches afin de mieux illustrer les cas reels de support, de supervision et d'escalade.

\section{Lecon generale}

La plupart de ces corrections ne changent pas l'idee du produit, mais elles changent fortement la qualite percue. Cela montre qu'un bon projet se joue aussi dans les details de coherence.

\chapter{Glossaire simplifie}

\begingroup
\small
\setlength{\tabcolsep}{4pt}
\begin{longtable}{>{\raggedright\arraybackslash}p{3cm} >{\raggedright\arraybackslash}p{9.8cm}}
\toprule
Terme & Definition \\
\midrule
Ticket & Dossier representant une demande de support ou un incident. \\
SLA & Engagement de delai de traitement ou de resolution. \\
Escalade & Reorientation ou supervision renforcee d'un ticket. \\
Camunda & Moteur BPMN charge d'executer le workflow metier. \\
Keycloak & Serveur d'identite et d'authentification. \\
Alfresco & GED utilisee pour l'archivage documentaire. \\
Ollama & Runtime local pour les modeles de langage. \\
FastAPI & Framework Python utilise pour l'agent IA. \\
RBAC & Controle d'acces base sur les roles. \\
Workflow & Enchainement structure des etapes du traitement. \\
Cockpit & Interface Camunda de supervision des processus. \\
Shortlist & Liste reduite de candidats proposes pour une assignation. \\
\bottomrule
\end{longtable}
\endgroup

